% CV generated from people/kevin/facts.yml according to brief.md.
% Template: templates/classic (see TEMPLATE.md).
\documentclass[10pt, a4paper]{article}
\usepackage[french]{babel}
\usepackage{kcv-classic}

\begin{document}

\cvheader{photo.jpg}{Kevin Bosirany ORLANDO}{%
  \cvlocation{Montpellier, France}\cvsep
  \cvmail{kevinbosirany@gmail.com}\cvsep
  \cvphone{+33762042840}{+33 (0)7 62 04 28 40}\cvsep
  \cvlink{LinkedinIn}{https://linkedin.com/in/kbosirany}{kbosirany}\cvsep
  \cvlink{Gitlab}{https://gitlab.com/kbosirany}{kbosirany}\cvsep
  \cvlink{Orcid}{https://orcid.org/0009-0009-2784-3108}{0009-0009-2784-3108}%
}

\section{Profil}
\cvtext{Ingénieur en modélisation hydrologique et agronomique avec deux ans
d'expérience en modélisation et étude d'impact du changement climatique sur les
ressources en eau. Compétences en modélisation des systèmes hydrologiques
complexes, en développement de modèles sous R et en traitement de données.
En quête d'un nouveau défi au sein d'un collectif où valoriser et développer
mes compétences dans un cadre opérationnel.}

\section{Expériences professionnelles}
\cventry{Ingénieur d'études}{Chargé de modélisation hydrologique et agronomique}%
  {INRAE -- UMR G-EAU / UMR Innovation}{Montpellier, France}{Sept 2023 -- Aujourd'hui}
\begin{cvpoints}
  \item Étude d'impact du changement climatique sur les volumes prélevables du
  bassin versant de l'Aude (projet Talanoa Water)
  \begin{itemize}
    \item Modélisation hydrologique semi-distribuée du réseau hydrographique et
    hydraulique du bassin versant (GR4J + CemaNeige)
    \item Acquisition et traitement des données climatiques (SAFRAN, DRIAS
    Explore 2), des prélèvements en eau (AERMC, DDTM…) et de la gestion des
    infrastructures hydrauliques (EDF, VNF, BRL…)
    \item Calibration du modèle et naturalisation des débits
    \item Calcul des volumes prélevables
  \end{itemize}
  \item Développement du modèle Optirrig (génération, analyse et optimisation de
  scénarios d'irrigation)
  \begin{itemize}
    \item Développement et maintenance du package R du modèle (OptirrigCORE)
    \item Développement et implémentation des formalismes du modèle (croissance
    de la plante, bilan hydrique…)
    \item Acquisition et traitement des données expérimentales
    \item Calibration du modèle pour différentes cultures (maïs, blé…)
  \end{itemize}
  \item Création d'une base de données spatiale des besoins en eau des
  principales cultures irriguées en France sur la période 2000--2100 (projet
  MAA -- PTGE)
\end{cvpoints}

\section{Stages}
\cventry{Stage de fin d'études}{}{INRAE -- UMR G-EAU / UMR Innovation}%
  {Montpellier, France}{Mars -- Août 2023}
\begin{cvpoints}
  \item Modélisation de la gestion de la ressource en eau du bassin versant de
  l'Aude dans un contexte de changement climatique (projet Talanoa Water)
\end{cvpoints}
\cventry{Stage de fin d'études}{}%
  {Institut Agronomique et Vétérinaire Hassan II -- Département Génie Rural}%
  {Rabat, Maroc}{Fév -- Juil 2022}
\begin{cvpoints}
  \item Détermination de la consommation et de la productivité de l'eau dans le
  bassin hydraulique de l'Oum Er Rbia
\end{cvpoints}
\cventry{Stage d'insertion en milieu professionnel}{}%
  {Compagnie d'Aménagement et de Développement Industriel (ADI)}%
  {Rabat, Maroc}{Juil -- Sept 2021}
\begin{cvpoints}
  \item Dimensionnement du réseau projeté pour la reconversion en irrigation
  localisée du secteur RD1 du périmètre de la Tessaout aval
\end{cvpoints}

\section{Parcours académique}
\cventry{Master 2 Sciences de l'Eau -- Eau et Agriculture}{}%
  {Université de Montpellier -- Institut Agro Montpellier -- AgroParisTech}%
  {Montpellier, France}{Sept 2022 -- Sept 2023}
\cventry{Ingénieur Génie Rural -- Eau, Environnement et Infrastructures}{}%
  {Institut Agronomique et Vétérinaire Hassan II}%
  {Rabat, Maroc}{Sept 2019 -- Juil 2022}
\cventry{Licence Génie Civil -- Hydraulique et Aménagement}{}%
  {École Supérieure Polytechnique d'Antananarivo}%
  {Antananarivo, Madagascar}{Sept 2015 -- Juil 2019}

\section{Compétences}
\cvskill{Programmation}{R, Python (bases), Fortran (bases)}
\cvskill{Modélisation}{GR4J, CemaNeige, Optirrig, Epanet}
\cvskill{Outils}{Git, GitLab, GitHub, CI/CD, Windows, Linux, Microsoft Office, LibreOffice, AutoCAD}

\end{document}
